% !Mode:: "TeX:UTF-8"

% 中英文摘要
\begin{cabstract}
随着移动互联网与社交媒体的蓬勃发展，包含文本、图像等多模态数据的用户生成内容已成为社会舆情表达的主要载体。多模态方面级情感分析(Multimodal Aspect-Based Sentiment Analysis, MABSA)旨在从多模态数据中识别针对特定评价对象的细粒度情感倾向，是当前情感分析领域的研究热点。然而，现有研究多侧重于图文间的显式语义匹配，往往忽视了图文之间隐含的情感关联。在复杂的社交媒体环境下，图像与文本的关系可能呈现互补、无关甚至矛盾的状态，这种情感指向的模糊性与噪声干扰导致模型在处理图文情感不一致场景时鲁棒性不足，且缺乏足够的可解释性。
针对上述挑战，本文从图文关系的视角出发，系统研究了基于图文关联的多模态方面级情感分析方法，旨在通过精准刻画图文交互逻辑来提升模型的情感理解能力。本文的主要研究内容与贡献如下：

首先，定义了基于情感表达的方面级图文相关性，并构建了相应的标注数据集。针对现有图文关系定义仅关注全局语义对齐而忽略局部情感关联的局限，本文提出了基于情感表达的细粒度相关性定义，并据此构建了方面级图文相关性数据集。通过系统的数据分析与基准实验，揭示了图文关系类型与情感极性判断之间的内在联系，验证了引入方面级跨模态相关性信息对提升模型性能的有效性。

其次，提出了一种基于图文情感线索引导的多模态方面级情感分析方法。针对现有方法对情感线索建模粒度粗糙的问题，该方法引入多模态大语言模型(Multimodal Large Language Model, MLLM)从图像和文本中提取与特定方面相关的情感线索描述，将其作为补充语义编码至上下文表示中。同时，利用细粒度方面级跨模态相关性作为引导信号，控制情感线索与图文特征的自适应融合。实验结果表明，该机制能够有效抑制无关模态或冲突线索的噪声干扰，显著增强了模型对复杂情感表达的理解能力。

最后，提出了一种基于多角度图文相关性解耦的多模态方面级情感分析方法。为进一步提升模型在复杂场景下的鲁棒性，本文将方面级跨模态相关性解耦为语义相关性与情感相关性两个独立维度。通过分别提取图像中的语义与情感特征，并利用解耦后的多角度相关性概率作为门控信号，实现了对特征融合的精细化控制。此外，构建了多任务学习机制，将单模态情感线索从输入端的提示信息转化为训练阶段的监督信号，促使模型主动捕获关键情感特征，从而在有效过滤视觉噪声的同时提升了情感分析的准确性。
	\vskip 21bp
	{\heiti\zihao{-4} 关键词：}
	多模态情感分析；
	方面级情感分析；
	图文关系；
	多模态大语言模型；
	
	\begin{flushright}
		作~~~~~~~~者：李兆煜
		
		指导老师：付国宏
		
	\end{flushright}
\end{cabstract}


